\clearpage\thispagestyle{otFront}\phantomsection\label{avant-propos}\pdfbookmark[0]{Apprendre et comprendre}{bm-avant-propos}\otPageTitre{Apprendre et comprendre}{Un livre pour toute la Quatrième}\fontsize{11}{15}\selectfont

Ce livre accompagne une année de mathématiques en Quatrième : consolider
les calculs, découvrir les outils géométriques et résoudre des problèmes
qui demandent de relier plusieurs idées. Il peut être suivi dans l’ordre
ou utilisé pour reprendre une difficulté précise.\par\medskip
\begin{savoirfaire}[Une progression régulière]Chaque chapitre contient des objectifs, des prérequis, deux pages de cours et quinze exercices : cinq pour découvrir, cinq pour s’entraîner et cinq problèmes. Les corrections expliquent le choix de méthode, les étapes et l’interprétation du résultat.\end{savoirfaire}
\begin{savoirfaire}[Des représentations pour penser]Les figures de géométrie, tableaux, graphiques, patrons et blocs de programmation accompagnent les questions qui en ont besoin. Les constructions cherchées et les réponses graphiques sont ajoutées au corrigé. Le dessin sert à comprendre et à chercher ; les données restent la base des preuves.\end{savoirfaire}
\begin{savoirfaire}[Chercher davantage]Dix-huit défis facultatifs proposent une recherche plus longue, avec deux indications séparées et une solution. Deux bilans annuels permettent de mélanger les notions. Les fiches outils et le repérage final orientent les révisions.\end{savoirfaire}
Les six compétences sont travaillées : chercher, modéliser, représenter,
raisonner, calculer et communiquer. Le progrès se mesure aussi à la capacité
d’expliquer une démarche et de corriger une erreur.

\clearpage
